\documentclass[12pt,a4paper]{article}
\usepackage{amsmath}
\usepackage{amssymb}
\usepackage{bm}
\usepackage{color}
\usepackage[utf-8]{inputenc}
\usepackage{xeCJK}

\setCJKmainfont{Hiragino Sans GB}

\title{京都大学大学院 数理工学専攻}
\author{凸最適化 過去問}
\date{H28 (OR)}

\begin{document}

\maketitle

\section*{問題 (i)}

不等式を不等号に。

\begin{enumerate}
  \item[\textcircled{1}] $x^*$ が定行可能解ですか。最適解 $\alpha$ はあり，$f(x^*) \leq f(x)$
  
  \item[\textcircled{2}] $f$ が凸性から 
  \[f(\alpha x + (1-t)y) \leq \alpha f(x) + (1-t)f(y)\]
  
  \item[\textcircled{3}] 以前の小問の不等式に適当な値を代入
\end{enumerate}

式と不等式を解で判断する。令回は $0$

\[(x^* - x^k)^T(x^T - x^k) \leq I \text{ より } x^k \text{ は } P(k) \text{ の実行可能解である。}\]

したがって，$f_k(x^*) \leq f_k(x^k)$

令 $0$ と $x^* = 0$ $f$ で考え $f_k(x^*) \leq f(x^*)$ より $f_k(x^*) \leq f(x^*)$

\section*{問題 (ii)}

\[\lim_k k(\omega x^k) = \bar{\lambda} \text{ より } \lim_k k(\rho x^k) \text{ か考え有限の他入に与えるが}\]

\[\omega t x^k = \frac{\bar{\lambda}}{k} + o\left(\frac{1}{k}\right) \text{ が必要である（$o$ はランダウの記号 - ）}\]

したがって，
\[\lim \partial \omega x^k = \partial \bar{x} = 0\]

手札：$\partial \omega \bar{x} = 0$ より $\bar{x}$ は $P$ の実行可能解である。

(ii)より 
\[\lim f_u(x^k) = f(\bar{x}) + \frac{1}{2}(\bar{x} - x^k)^T(\bar{x} - x^k) \leq f(x^k)\]

もし $\bar{x} \neq x^k$ ならば $(\bar{x} - x^k)^T(\bar{x} - x^k) > 0$ ゆえ $x^k$ が $P$ の大局的最適解である
ここに矛盾 $f$。したがって，$\bar{x} = x^k$

\section*{問題 (iii)}

$\lambda$ をうがランジュ乗数とするとき，$P(k)$ のKKT条件は
\begin{align}
\begin{cases}
\nabla f_k(x) + 2\lambda(x - x^k) = 0 \\
(x - x^k)^T(x - x^k) - I \leq 0, \quad \lambda \geq 0, \quad \lambda((x - x^k)^T(x - x^k) - 1) = 0
\end{cases}
\tag{★}
\end{align}

\section*{問題 (iv)}

$x^*$ はKKT条件と誤差が $\nabla f(x^k) + 2\lambda(x^k - x^k) = 0$，$\lambda = 0$ を示したい。

$\rightarrow$ 相対性優 KKT条件 $\lambda((x^k - x^k)^T(x^k - x^k) - 1) = 0$ より $\lambda = 0$ を示す。

$x^k$ は $P(k)$ の最適解が KKT条件と満たすので入が合致して
\[\nabla f_k(x^k) + 2\lambda(x^k - x^k) = 0\]

令(ii)より $x^k \to x^k$（$k \to \infty$）をかく十分大きく に対して 
\[(x^k - x^k)^T(x^k - x^k) < 1\]

となる。この時，$\lambda((x^k - x^k)^T(x^k - x^k) - 1) = 0$ より $\lambda = 0$

したがって，$2$，$k$ が十分大きいとき，$\nabla f_k(x^k) = 0$

\section*{問題 (v)}

(iv)より 十分大きな$k$に対して $\nabla f_k(x^k) = 0$，より
\[\nabla f(x^k) + k \partial \omega x^k \partial \omega + (x^k - x^k) = 0\]

$k \to \infty$の和限度として(ii)より $\lim_{k \to \infty} x^k = x^k$である

\[\nabla f(x^k) + \bar{\lambda} \partial \omega = 0\]

\newpage

\section*{Remind}

\subsection*{ランダウの記号}

\[f(x) = O(x^2) + 2x^c\]

の如く，$x$が十分大きいと。その挙動は $O(x^2)$と呼称する。こんなランダウの\textbf{ビッグオー} - $\mathcal{O}$ を用いて示す。

\[f(x) = O(\omega^2) \quad (x \to \infty)\]

また，$f(x)$は$x$が十分大きいと，$\omega x^k$よりも$3$分の$1$小さいのこれを
ランダウの\textbf{小記号} - $o$ と用いて示す。

\[f(x) = o(x^4) \quad (x \to \infty)\]

逆に対$0$に連いては
\[f(x) = O(x) \quad (x \to 0), \quad f(x) = o(x) \quad (x \to 0)\]

\end{document}
